\section{Resultados e discussão}\label{sec:results}

\subsection{Qual verificação governa, e onde}\label{sec:governante}

A premissa deste estudo é que a verificação governante migra com o vão, e que a propriedade do material que determina o dimensionamento migra com ela. Essa premissa é verificada em primeiro lugar, porque tudo o que se segue depende dela.

\begin{table}[!ht]
\caption{Verificação governante da solução de volume mínimo, por vão e classe de resistência.\label{tab:governante}}
\begin{threeparttable}
\begin{tabular*}{\columnwidth}{@{\extracolsep\fill}lccccc@{\extracolsep\fill}}
\toprule
 & \multicolumn{5}{c}{Classe de resistência} \\
\cmidrule{2-6}
Vão $L$ (m) & D20 & D30 & D40 & D50 & D60 \\
\midrule
3,0 & \tofill{--} & \tofill{--} & \tofill{--} & \tofill{--} & \tofill{--} \\
5,0 & \tofill{--} & \tofill{--} & \tofill{--} & \tofill{--} & \tofill{--} \\
8,0 & \tofill{--} & \tofill{--} & \tofill{--} & \tofill{--} & \tofill{--} \\
10,0 & \tofill{--} & \tofill{--} & \tofill{--} & \tofill{--} & \tofill{--} \\
\bottomrule
\end{tabular*}
\begin{tablenotes}
\footnotesize
\item Legenda: C = cisalhamento da longarina; F = flexão da longarina; D = deslocamento da longarina; FT = flexão da peça do tabuleiro. A verificação governante é aquela com a maior função de estado limite normalizada na solução de volume mínimo.
\end{tablenotes}
\end{threeparttable}
\end{table}

\begin{fillblock}
\noindent\textbf{[A PREENCHER -- resultado que sustenta toda a tese]} Preencher a Tabela~\ref{tab:governante} e verificar se a migração esperada de fato ocorre: cisalhamento ou flexão governando nos vãos curtos e flecha governando nos vãos longos. Reportar em que vão, aproximadamente, ocorre a transição, e se ela depende da classe de resistência.

Este resultado precisa ser reportado com honestidade em duas direções. Se a migração ocorrer conforme esperado, ela é o mecanismo do artigo e deve ser declarada explicitamente aqui, com a consequência de que o erro de enquadramento de rigidez só se manifesta estruturalmente na região em que a flecha governa. Se a flecha governar em toda a faixa de vãos, o argumento continua válido -- e até mais forte, pois significa que a rigidez sempre importa -- porém a dependência do vão deixa de ser o eixo do artigo e o texto deve ser reescrito em torno da magnitude do erro, não da sua variação. Se, ao contrário, a flecha nunca governar, o efeito estrutural será pequeno e isso precisa ser dito: nesse caso o achado do artigo passa a ser que o erro de representação de rigidez existe mas tem consequência limitada nesta tipologia, o que é um resultado negativo publicável, mas exige reescrever a introdução e o resumo.
\end{fillblock}

\subsection{Custo do enquadramento em classe}\label{sec:custo_enquadramento}

Para cada espécie e vão, a Tabela~\ref{tab:custo_enquadramento} reporta o volume de madeira do projeto pela espécie, o do projeto pela classe e o erro $\Delta V$ resultante, definido na Equação~\eqref{eq:deltaV}.

\begin{table}[!ht]
\caption{Volume de madeira dos projetos pela espécie e pela classe e o erro de enquadramento resultante, por espécie.\label{tab:custo_enquadramento}}
\begin{threeparttable}
\footnotesize
\setlength{\tabcolsep}{3pt}
\begin{tabular*}{\textwidth}{@{\extracolsep{\fill}}llcrrrrrr@{\extracolsep{\fill}}}
\toprule
\# & Espécie & Classe & \multicolumn{2}{c}{$L = 3{,}0$~m} & \multicolumn{2}{c}{$L = 5{,}0$~m} & \multicolumn{2}{c}{$L = 10{,}0$~m} \\
\cmidrule(lr){4-5}\cmidrule(lr){6-7}\cmidrule(lr){8-9}
 & & & $V_{sp}$ & $\Delta V$ & $V_{sp}$ & $\Delta V$ & $V_{sp}$ & $\Delta V$ \\
\midrule
\multicolumn{9}{c}{\tofill{[A PREENCHER -- 40 linhas, uma por espécie, na mesma ordem da Tabela~\ref{tab:especies}]}} \\
\midrule
\multicolumn{3}{l}{$|\Delta V|$ médio} & & \tofill{--\%} & & \tofill{--\%} & & \tofill{--\%} \\
\multicolumn{3}{l}{Amplitude de $\Delta V$} & & \tofill{--\%} & & \tofill{--\%} & & \tofill{--\%} \\
\multicolumn{3}{l}{Espécies com $\Delta V < 0$} & & \tofill{--} & & \tofill{--} & & \tofill{--} \\
\bottomrule
\end{tabular*}
\begin{tablenotes}
\footnotesize
\item $V_{sp}$: volume de madeira do projeto pela espécie (m\textsuperscript{3}). $\Delta V$: erro de enquadramento em classe, Equação~\eqref{eq:deltaV}. Um $\Delta V$ negativo significa que o projeto pela classe emprega menos material do que a espécie requer, sendo portanto contrário à segurança.
\item \tofill{[NOTA] Caso a tabela completa de 40 linhas por três vãos fique excessiva para o corpo do artigo, mover a versão completa para material suplementar e manter aqui apenas as dez espécies de maior $|\Delta V|$ mais as linhas de síntese. Verificar a política de material suplementar da revista.}
\end{tablenotes}
\end{threeparttable}
\end{table}

\begin{figure}[!ht]
    \centering
    \includegraphics[width=0.9\textwidth]{figuras/continuo_vs_degraus.png}
    \caption{Volume de madeira como função contínua da propriedade do material contra a função-degrau imposta pelo projeto baseado em classe, um painel por vão. Cada ponto é uma espécie dimensionada com suas propriedades medidas; a função-degrau é o projeto pela classe.}
    \label{fig:continuo_degraus}
\end{figure}

\begin{fillblock}
\noindent\textbf{[A PREENCHER -- figura central do artigo]} Construir a Figura~\ref{fig:continuo_degraus} com um painel por vão, dispostos lado a lado e compartilhando o eixo vertical. Em cada painel: no eixo horizontal a resistência característica à compressão, no eixo vertical o volume de madeira do projeto; um ponto por espécie, obtido com as propriedades medidas; e sobreposta, em traço escalonado, a função-degrau correspondente ao projeto feito pela classe. A distância vertical entre cada ponto e o degrau da sua faixa é o erro de enquadramento daquela espécie.

A leitura pretendida é que a nuvem de pontos se espalhe verticalmente em torno do degrau, e que esse espalhamento cresça de painel para painel conforme o vão aumenta. Se for esse o caso, a figura demonstra sozinha a tese do artigo e deve ser referenciada no resumo. Marcar nominalmente, em todos os painéis, as duas espécies extremas discutidas na Seção~\ref{sec:representacao}, para que o leitor possa acompanhá-las.
\end{fillblock}

\begin{fillblock}
\noindent\textbf{[A PREENCHER -- discussão]} Discutir a Tabela~\ref{tab:custo_enquadramento} e a Figura~\ref{fig:continuo_degraus} em três frentes.

\textbf{(i) Magnitude.} Reportar o erro médio absoluto e a amplitude de $\Delta V$ em cada vão. Comparar com o erro de representação de rigidez da Seção~\ref{sec:representacao}: espera-se que o erro em volume seja da mesma ordem, ou amplificado, dependendo de quão diretamente a rigidez entra na restrição ativa.

\textbf{(ii) Direção.} Quantificar quantas espécies apresentam $\Delta V$ negativo, isto é, quantas recebem um projeto mais leve do que suas propriedades reais comportam. Este é o achado de segurança e deve ser reportado sem atenuação. Relacionar a lista dessas espécies com as 14 identificadas na Tabela~\ref{tab:discordancia} como menos rígidas que a classe: as duas listas devem coincidir em larga medida, e a coincidência é uma verificação de consistência interna do estudo.

\textbf{(iii) Dependência do vão.} Verificar se o erro cresce com o vão, conforme previsto pela migração do critério governante da Seção~\ref{sec:governante}. Quantificar o crescimento. Esta é a afirmação central do artigo e precisa vir acompanhada de número, não apenas de tendência.
\end{fillblock}

\subsection{A limitação está no número de classes?}\label{sec:granularidade}

Uma objeção natural aos resultados precedentes é que eles refletiriam a grosseria de um sistema de cinco classes, e seriam mitigados por um sistema mais fino. A EN 338 define treze classes de folhosas, de D18 a D80, e oferece portanto um teste direto dessa hipótese sobre a mesma população.

\begin{table}[!ht]
\caption{Comparação do erro de enquadramento sob os sistemas de classes da ABNT NBR 7190-3 e da EN 338.\label{tab:comparacao_normas}}
\begin{threeparttable}
\begin{tabular*}{\columnwidth}{@{\extracolsep\fill}lcc@{\extracolsep\fill}}
\toprule
 & ABNT NBR 7190-3 & EN 338 \\
\midrule
Número de classes de folhosas & 5 & \tofill{13} \\
Propriedade indexadora & $f_{c0,k}$ & \tofill{$f_{m,k}$} \\
Passo em $E$ tabelado entre classes vizinhas & 14--21\% & \tofill{--} \\
Dispersão intraclasse média de $E_{M0}$ medido & \tofill{--} & \tofill{--} \\
Razão dispersão intraclasse / passo entre classes & \tofill{$\approx 4$} & \tofill{--} \\
Espécies enquadradas em outra classe pela rigidez & 58\% & \tofill{--} \\
Espécies menos rígidas que sua classe & 35\% & \tofill{--} \\
\bottomrule
\end{tabular*}
\end{threeparttable}
\end{table}

\begin{fillblock}
\noindent\textbf{[A PREENCHER -- custa uma planilha, vale a internacionalização do artigo]} Reenquadrar as mesmas 40 espécies nas classes da EN 338 e preencher a Tabela~\ref{tab:comparacao_normas}. \emph{Nenhuma alteração no modelo de cálculo é necessária}: trata-se apenas de uma segunda regra de atribuição de classe aplicada à mesma população.

Dois cuidados. Primeiro, conferir os valores da EN 338:2016, Tabela 2, contra o texto normativo -- não reproduzir de memória nem de fonte secundária. Segundo, atentar para o fato de que a EN 338 indexa as classes de folhosas pela resistência à flexão, e não pela compressão como a NBR; isso não é um detalhe, e sim parte do argumento, pois permite testar se a escolha da propriedade indexadora importa mais que a granularidade.

O desfecho esperado, e que deve ser declarado com clareza qualquer que seja: se a EN 338, com treze classes, ainda assim apresentar dispersão intraclasse superior ao passo entre classes, então a limitação não está no número de classes e sim em indexar o sistema por resistência -- e essa é a conclusão geral que torna o artigo relevante fora do Brasil. Se, ao contrário, a granularidade mais fina resolver o problema, a conclusão passa a ser uma recomendação concreta de revisão da NBR 7190-3, o que é um resultado menos geral mas de aplicação imediata.
\end{fillblock}

\subsection{Caso detalhado}\label{sec:caso_detalhado}

\begin{fillblock}
\noindent\textbf{[A PREENCHER]} Selecionar uma espécie de referência de relevância comercial e apresentar, para ela, a fronteira eficiente completa, a análise de sensibilidade global de Sobol e a curva de convergência do hipervolume, de modo a documentar a qualidade da otimização que sustenta todos os resultados agregados das seções anteriores. Acrescentar a comparação resumida entre os três extremos do conjunto: a espécie mais eficiente, a menos eficiente e a de maior discrepância frente à sua classe.

Manter esta seção contida. Ela existe para dar credibilidade ao procedimento de otimização, não para reapresentar a plataforma.
\end{fillblock}

\subsection{Limitações}\label{sec:limitacoes}

\begin{fillblock}
\noindent\textbf{[A PREENCHER]} Declarar as limitações de forma direta, pois todas serão levantadas na revisão e é preferível antecipá-las:

\begin{itemize}[nosep]
\item os dados provêm de uma única fonte experimental, e a conversão de valores médios em característicos segue a relação normativa, sem verificação independente por espécie;
\item a variabilidade dentro de cada espécie não foi considerada, apenas a variabilidade entre espécies; a robustez incorporada à otimização é de natureza geométrica, não material;
\item o sistema estrutural restringe-se à superestrutura de vão único simplesmente apoiado, sem ligações, encontros nem fundações;
\item o modelo estrutural é analítico, baseado em linhas de influência de viga biapoiada, sem análise por elementos finitos nem validação experimental;
\item não há comparação com uma ponte efetivamente construída.
\end{itemize}

Encerrar indicando que a incorporação da variabilidade material por meio de análise de confiabilidade, permitindo avaliar a uniformidade do índice de confiabilidade entre espécies de uma mesma classe, é o desdobramento natural e está em curso.
\end{fillblock}
